% =====================================================================
%  DRIFTGUARD — Section V: Operational Viability and Safety.
%  Narrative rewrite. The verifier is presented in a BALANCED way: it is a
%  genuinely good label-quality instrument; whether to use it to gate
%  retraining is a budget/horizon-dependent trade-off, not a flat negative.
%  Figures: fig_verifier_overview.pdf, fig_injection_live.pdf.
%  Verifier ablation exact numbers + health-signal thresholds -> appendix.
% =====================================================================

\section{Operational Viability and Safety}
\label{sec:safety}

Recovering drift is one thing; running the loop safely is another. In
deployment the oracle is an LLM, its input is attacker-controlled page content,
and there is no ground truth at runtime to check it. This section builds the
safety layer that makes the loop deployable anyway, and is candid about its
limits. A class-asymmetric verifier scores every oracle label without an LLM
call; we show it is a good label-quality instrument, then ask -- and answer
with nuance -- whether it should \emph{gate} retraining. In addition, a label-free
health monitor is the loop's single circuit breaker, validated on oracles that
failed for real --- and we are explicit about the one failure mode it provably
cannot see, which we carry forward to future work.

\subsection{Scoring Oracle Labels without Ground Truth}
\label{ssec:verifier}

The oracle's two verdict classes are not equally reliable: on the AL stream a
phishing verdict is $\approx\!99\%$ precise, but a benign verdict only
$\approx\!59\%$ --- clean-looking scam pages the LLM excuses. The verifier is
built around that asymmetry. For each label it produces a soft trust score in
$[0,1]$ from three tiers: the model's own confidence, an evidence check that
dispatches each claimed indicator to one of a dozen HTML/URL/OSINT predicates
and asks whether it is actually present, and an OSINT-coherence tier that
\emph{reinforces} a phishing verdict but \emph{challenges} a benign one. For example a benign
call on a domain registered days ago, never archived and absent from Tranco is
penalised. The exact predicates and weights are in the appendix; two properties
matter here. It never needs an LLM call, and --- by design --- it never discards a
label on its own: it only measures how well a verdict is grounded in checkable
evidence.

As a measurement, it works well (Figure~\ref{fig:verifier}). Restricting the
canonical $1{,}247$-label stream to the labels the verifier trusts
($\tau\!\geq\!0.5$, $n{=}1{,}135$) raises oracle accuracy from $0.83$ to $0.89$
and, on the weak phishing class the loop most needs, oracle recall from $0.57$ to
$0.68$. The $112$ labels it distrusts are $72\%$ wrong, and essentially all of
them are the same failure: a genuinely phishing page the oracle excused as
benign. The verifier is, in short, a faithful label-quality instrument. The
question is what to \emph{do} with that measurement.

\subsection{To Gate, or to Keep Every Label?}
\label{ssec:verifier-gate}

The obvious move is to act on the trust score at training time --- softly, as a
per-label weight, or hard, dropping every label below $0.5$. We test both against
the no-verifier baseline in a \emph{controlled} ablation that replays the
identical $1{,}247$-label trajectory, so the only variable is the weight a label
receives (Figure~\ref{fig:verifier}, lines).

The result is a genuine trade-off, and reading it correctly matters. In
distribution, cleaning the labels costs essentially nothing: the full verifier
lands at $\mathrm{PP}\,F_{1,w}=0.799$ and the hard filter at $0.804$, against the
$0.800$ baseline, with KP flat throughout --- so feeding the model a demonstrably
cleaner label stream neither helps nor hurts what it does on the drifted
distribution. The one place gating leaves a mark is out of distribution, and it
is the mark Section~\ref{sec:noise} predicts: the labels the verifier removes are
precisely the imperfect ones that were acting as a regularizer, so dropping them
sharpens the boundary back toward PP and costs the adversarial probe $4$--$8$
points of macro-$F_1$, the more aggressively the more it filters. A second,
independent instrument thus reproduces the cursor of Section~\ref{sec:noise}: the
``noisy'' labels are not defects to be filtered, they are the regularizer.

We therefore keep \emph{every} label at uniform weight in DRIFTGUARD, preserving
the out-of-distribution robustness of Section~\ref{sec:noise}. But this is a
decision about a \emph{regime}, not a verdict that label hygiene is worthless.
The reason hygiene buys nothing measurable here is that the $7{,}221$-sample seed
corpus dominates the retrain loss, so re-weighting a thousand oracle labels
barely moves the boundary; the regularizing noise, by contrast, is felt on the
sensitive adversarial axis. In a long-running deployment the balance shifts: as
acquired labels accumulate and eventually rival the seed, and as label errors
compound across many retrains, a verifier that lifts label accuracy from $0.83$
to $0.89$ becomes an asset. The honest reading
is that the verifier's value is real and grows with scale and horizon, while its
small-budget out-of-distribution cost is simply the flip side of the very noise
that helps today --- and that, at our budget, keeping the labels is the right call.

\begin{figure}[t]
  \centering
  \includegraphics[width=\columnwidth]{figures/fig_verifier_overview.pdf}
  \caption{The verifier works --- which is what makes the gating question
  interesting. Bars: across modes, restricting to the labels the verifier trusts
  raises oracle accuracy ($0.83\!\to\!0.89$) and phishing recall
  ($0.57\!\to\!0.68$). Lines: feeding that cleaner stream to the retrain leaves
  in-distribution PP flat (even marginally higher for the hard filter) but costs
  out-of-distribution robustness (DP), exactly the regularization of
  Section~\ref{sec:noise}. DRIFTGUARD keeps every label at uniform weight; the
  hygiene the verifier offers is a budget- and horizon-dependent asset, not a
  one-shot $F_1$ boost.}
  \label{fig:verifier}
\end{figure}

\subsection{The Verifier as a Cross-Model Auditor}
\label{ssec:auditor}

If the verifier should not gate individual labels, its real, regime-independent
job is \emph{observability}. Because it scores the \emph{grounding} of a verdict
rather than its outcome, and does so without ground truth, it yields per-model
behavioural metrics that accuracy alone cannot: how often a provider's claims are
refutable (a hallucination proxy), how well its confidence tracks correctness,
and what reasoning style it favours. This makes it the natural instrument for
choosing an oracle and watching it drift across models --- exactly how it is used
in Section~\ref{sec:oracle-quality} --- and it feeds the health monitor one
round-level trust aggregate, never a per-label decision. That benchmark also
closes the present argument from the other side: it surfaces two robust paradoxes
--- the verifier's highest-trust labels are \emph{less} accurate than its
second-highest, and the samples it flags as hallucinating are \emph{more}
accurate than clean ones --- which independently confirm that neither trust nor
hallucination is safe as a per-label gate.

\subsection{The Label-Free Health Monitor}
\label{ssec:safety-layer}

The verifier measures but never acts; the single actuator is a health monitor
that detects a structurally broken oracle from the output stream alone, with no
ground truth. An oracle can degrade silently --- a wrong model behind a gateway, a
rate-limited provider returning malformed JSON, a page attacking the prompt ---
and quietly poison the retrain. Each round the monitor computes a few label-free
signals over the returned verdicts: the \emph{parse rate} (schema-valid
fraction), the \emph{class skew}, the \emph{yield}, and optionally the verifier's
round-mean trust. The parse rate is the primary discriminator, and deliberately a
parse-layer rather than a network-layer check: a broken oracle can return
HTTP~200 with an unparseable body, whereas a healthy one essentially always emits
schema-conforming output, so its collapse has no benign explanation. Skew, by
contrast, is \emph{not} treated as breakage on its own --- a drift-anchored batch
is legitimately near-mono-class early in the stream --- so its critical threshold
is set at true mono-class and no single signal can halt the run. The monitor
stops only on corroborated, sustained failure (two consecutive multi-signal
critical rounds, or three consecutive rounds of degraded parse rate); it is a
pure observer that never edits a label, so it cannot bias learning.

The clearest evidence that this is the right design is a case where it fired for
real (Figure~\ref{fig:llama-halt}). Swapping the oracle for a
\textsc{Llama-3.3-70B} endpoint, the loop runs healthy for fifteen rounds and
then begins to degrade: the provider's schema-valid fraction erodes round by
round until three consecutive rounds fall below the parse-warning line, tripping
the sustained-parse rule and halting the loop at round~18. The decisive detail is
what the monitor does \emph{not} see. It reads only the parse-ok rate; it is
blind to the test-time $F_{1,w}$ of the model being retrained --- which, on this
oracle, is quietly sliding the whole time. Acting on the structural signal alone
is nonetheless enough: the halt stops the loop before the degrading labels can do
lasting harm, and the monitored run ends strictly above where the unmonitored run
would have driven it. The same parse-layer logic catches the two other ways the
oracle failed in our benchmark, without a single false alarm: a \textsc{Grok-4}
endpoint that was dead on arrival ($0/20$ usable labels), and a \textsc{Kimi-K2.6}
whose calls all returned HTTP~200 yet were $68\%$ unparseable --- which a network
liveness check would have waved through and the parse rate catches at once. Across
the four healthy providers it never tripped in $25$ rounds each; its full
cross-model validation is reported with the multi-LLM benchmark
(Section~\ref{sec:oracle-quality}).

\begin{figure}[t]
  \centering
  \includegraphics[width=\columnwidth]{figures/fig_llama_halt.pdf}
  \caption{The health monitor halting a real failing oracle. With a
  \textsc{Llama-3.3-70B} oracle the schema-valid (parse-ok) rate --- the only
  signal the monitor reads --- erodes after round~15 and trips the sustained-parse
  rule (three consecutive rounds below $0.95$), halting the loop at round~18. The
  monitor is \emph{blind} to the model's test $F_{1,w}$ (dotted, right axis),
  which degrades throughout; acting on the structural signal alone is enough to
  stop the loop before the broken oracle does lasting harm.}
  \label{fig:llama-halt}
\end{figure}

The honest limit follows directly from how the monitor works: it is a
\emph{structural-breakage} detector --- malformed output, sustained parse failure,
gross class collapse --- not a \emph{label-correctness} detector. An oracle that
stays perfectly schema-conforming while a prompt-injected page coerces it into a
confident, well-formed \emph{wrong} verdict produces no signal a threshold can
catch, and no threshold change closes that gap. The natural defence is the
verifier's OSINT-coherence tier, which by construction challenges a benign verdict
on a freshly registered, unarchived domain that a label-flip cannot also forge;
turning that observability into a safe runtime signal, without re-introducing the
gating pathology of Section~\ref{ssec:verifier-gate}, is a problem we take up
explicitly in Section~\ref{sec:conclusion}.

\medskip\noindent
The two components form the operational safety layer with a clean division of
labour: the verifier \emph{measures} (per-label grounding and a round-level trust
signal, never gating), the monitor \emph{acts} (a sustained, corroborated halt,
never on a single label). Together they let the loop run without runtime ground
truth --- keeping every paid label, and with it the implicit-regularization
benefit of Section~\ref{sec:noise}, auditing the oracle continuously, and
stopping itself before a structurally broken oracle can poison the model. The
boundary of that guarantee is explicit: it defends against an oracle that
\emph{breaks}, not against one quietly coerced into confident, well-formed lies
--- the open problem we carry to Section~\ref{sec:conclusion}.
